\documentclass[12pt]{article}
\usepackage[utf8]{inputenc}
\usepackage[T1]{fontenc}
\usepackage[margin=1in]{geometry}
\usepackage{setspace}
\usepackage{xcolor}
\usepackage[colorlinks=true, urlcolor=blue!60!black, linkcolor=blue!60!black]{hyperref}
\setlength{\parindent}{2em}
\setlength{\parskip}{1.25em}
\renewcommand{\baselinestretch}{1.0}
\begin{document}
\begin{center}
{\LARGE \textit{United States v. Virginia}, 518 U.S. 515 (1996)\par}
\vspace{0.8em}
{\large Case Brief\par}
\end{center}
\vspace{1em}

\section{Memory Jogger}
In \href{https://www.courtlistener.com/opinion/118056/united-states-v-virginia/}{\textit{United States v. Virginia}}, 518 U.S. 515 (1996), the Supreme Court held 7--1 that Virginia violated the Fourteenth Amendment's Equal Protection Clause by excluding women from the Virginia Military Institute without an ``exceedingly persuasive justification'' and that creating a separate, inferior women's leadership program did not cure the constitutional violation.

\section{Detailed Case Facts}
Founded in 1839 under state statute (1839 Va. Acts, ch. 20) and remaining financially supported by the Commonwealth and subject to the control of the Virginia General Assembly (\href{https://supreme.justia.com/cases/federal/us/518/515/}{Va. Code Ann. \S~23-92 (1993)}), the Virginia Military Institute (VMI) was the sole single-sex institution among Virginia's fifteen public colleges and universities (\href{https://supreme.justia.com/cases/federal/us/518/515/}{518 U.S. at 520--521}). Enrolling approximately 1,300 male cadets, VMI pursued a distinctive mission of producing ``citizen-soldiers'' prepared for leadership in both civilian life and military service, although only about 15\% of its cadets entered career military service (\href{https://supreme.justia.com/cases/federal/us/518/515/}{id. at 521--522}).

To accomplish this mission, VMI employed an ``adversative method'' of education characterized by physical rigor, mental stress, absolute equality of treatment, absence of personal privacy in spartan barracks, minute regulation of behavior, a seven-month ``rat line'' for entering cadets comparable in intensity to Marine Corps boot camp, a hierarchical class and ``dyke'' mentoring system, and a strictly enforced honor code (\href{https://supreme.justia.com/cases/federal/us/518/515/}{id. at 522}). VMI's loyal alumni network included military generals, Members of Congress, and business executives, and the institute boasted the largest per-student endowment of any public undergraduate institution in the United States---\$131 million at the time of trial, with \$220 million in future commitments (\href{https://supreme.justia.com/cases/federal/us/518/515/}{id. at 520, 527, 552}).

In the two years preceding the litigation, VMI received admission inquiries from 347 women and responded to none of them (\href{https://supreme.justia.com/cases/federal/us/518/515/}{id. at 523}). In 1990, prompted by a complaint filed with the United States Attorney General by a female high-school student seeking admission to VMI, the United States sued the Commonwealth of Virginia and VMI, alleging that VMI's male-only admissions policy violated the Equal Protection Clause of the Fourteenth Amendment (\href{https://supreme.justia.com/cases/federal/us/518/515/}{id.}). The VMI Foundation and the VMI Alumni Association intervened as defendants (\href{https://supreme.justia.com/cases/federal/us/518/515/}{id. at 523 n.3}).

Following an initial appellate ruling that VMI's exclusionary policy was unconstitutional, Virginia proposed a remedial alternative: the Virginia Women's Institute for Leadership (VWIL), a four-year, state-sponsored program housed at Mary Baldwin College, a private liberal arts women's college, initially open to 25 to 30 students (\href{https://supreme.justia.com/cases/federal/us/518/515/}{id. at 526}). Designed by a Mary Baldwin Task Force around what it deemed appropriate for ``most women,'' VWIL replaced VMI's adversative military model with a ``cooperative method which reinforces self-esteem,'' omitted military barracks life and mandatory daily uniforms, and substituted leadership seminars, externships, community service, and a ``largely ceremonial'' Virginia Corps of Cadets (\href{https://supreme.justia.com/cases/federal/us/518/515/}{id. at 526--527, 548--549}). Tangible and intangible disparities were stark: Mary Baldwin freshmen averaged roughly 100 SAT points lower than VMI freshmen; Mary Baldwin's faculty held significantly fewer Ph.D.s and earned lower salaries; Mary Baldwin offered only Bachelor of Arts degrees and no on-campus engineering or advanced math and physics programs (requiring VWIL students seeking an engineering degree to attend Washington University in St. Louis for two years at private tuition); its athletic facilities consisted of two multipurpose fields and one gymnasium compared to VMI's extensive NCAA-level complex; and Mary Baldwin's \$19 million endowment (plus \$35 million in commitments and a \$5.4625 million VMI Foundation grant) paled beside VMI's financial reserves, 157-year prestige, and alumni network (\href{https://supreme.justia.com/cases/federal/us/518/515/}{id. at 526--527, 551--553}).

\section{Procedural History}
\begin{itemize}
\item \textbf{U.S. District Court (Phase I):} After a six-day bench trial, the U.S. District Court for the Western District of Virginia entered judgment for Virginia and VMI in \textit{United States v. Virginia}, 766 F. Supp. 1407 (W.D. Va. 1991), holding that single-sex education added diversity to Virginia's higher-education system and that admitting women would materially alter VMI's adversative method (\href{https://supreme.justia.com/cases/federal/us/518/515/}{518 U.S. at 523--524}).
\item \textbf{U.S. Court of Appeals (Phase I):} The Fourth Circuit vacated and remanded in \textit{United States v. Virginia}, 976 F.2d 890 (4th Cir. 1992), holding that Virginia had not justified affording VMI's unique program exclusively to men under an asserted diversity policy, and offering Virginia three remedial paths: admit women to VMI, create a parallel program for women, or privatize VMI (\href{https://supreme.justia.com/cases/federal/us/518/515/}{id. at 524--526}).
\item \textbf{Supreme Court (Interlocutory):} The Supreme Court denied Virginia's petition for certiorari in \textit{Virginia v. United States}, 508 U.S. 946 (1993), with Justice Scalia writing separately to note the interlocutory posture of the case (\href{https://supreme.justia.com/cases/federal/us/518/515/}{id. at 526}).
\item \textbf{U.S. District Court (Phase II):} On remand, Virginia proposed VWIL, and the District Court approved the remedial plan in \textit{United States v. Virginia}, 852 F. Supp. 471 (W.D. Va. 1994), concluding that VMI and VWIL would ``achieve substantially similar outcomes'' (\href{https://supreme.justia.com/cases/federal/us/518/515/}{id. at 527--528}).
\item \textbf{U.S. Court of Appeals (Phase II \& En Banc):} A divided Fourth Circuit panel affirmed in \textit{United States v. Virginia}, 44 F.3d 1229 (4th Cir. 1995), reviewing Virginia's objective deferentially and applying a newly invented ``substantive comparability'' test; Senior Judge Phillips dissented, calling VWIL a ``pale shadow'' of VMI (\href{https://supreme.justia.com/cases/federal/us/518/515/}{id. at 528--530}). The Fourth Circuit denied rehearing en banc in 52 F.3d 90 (4th Cir. 1995) (6--4 with three recusals, falling short of a majority of the thirteen active judges), over a dissent by Judge Motz joined by Judges Hall, Murnaghan, and Michael (\href{https://supreme.justia.com/cases/federal/us/518/515/}{id. at 530 \& n.4}).
\item \textbf{U.S. Supreme Court:} Both parties sought review (No. 94-1941 and No. 94-2107), and the Court granted certiorari (516 U.S. 910 (1995)), heard argument on January 17, 1996, and issued its decision on June 26, 1996 (\href{https://www.oyez.org/cases/1995/94-1941}{Oyez}).
\end{itemize}

\section{Judicial Votes}
The Supreme Court decided the case by a \textbf{7--1} vote, with one recusal (\href{https://supreme.justia.com/cases/federal/us/518/515/}{518 U.S. at 518, 558}):
\begin{itemize}
\item \textbf{Majority (6):} Justice Ruth Bader Ginsburg delivered the opinion of the Court, joined by Justices John Paul Stevens, Sandra Day O'Connor, Anthony M. Kennedy, David H. Souter, and Stephen G. Breyer.
\item \textbf{Concurring in the Judgment (1):} Chief Justice William H. Rehnquist filed an opinion concurring in the judgment.
\item \textbf{Dissenting (1):} Justice Antonin Scalia filed a dissenting opinion.
\item \textbf{Not Participating (1):} Justice Clarence Thomas took no part in the consideration or decision of the case (his son was enrolled at VMI at the time, \href{https://en.wikipedia.org/wiki/United_States_v._Virginia}{Wikipedia}).
\end{itemize}

\section{Holding}
\textbf{Judgment:} The Supreme Court \textbf{affirmed} the Fourth Circuit's initial liability judgment (976 F.2d 890) holding VMI's male-only admissions policy unconstitutional, \textbf{reversed} the Fourth Circuit's final remedial judgment (44 F.3d 1229) approving the VWIL program, and \textbf{remanded} for further proceedings consistent with its opinion (\href{https://supreme.justia.com/cases/federal/us/518/515/}{518 U.S. at 558}).

\textbf{Rule Announced:} Under the Fourteenth Amendment's Equal Protection Clause, a party seeking to defend government action classifying individuals by sex must demonstrate an ``exceedingly persuasive justification'' (\href{https://supreme.justia.com/cases/federal/us/518/515/}{id. at 531}). The burden rests entirely on the State, which must show at least that the classification serves important governmental objectives and that the discriminatory means employed are substantially related to achieving those objectives (\href{https://supreme.justia.com/cases/federal/us/518/515/}{id. at 533}). The justification must be genuine---not hypothesized or invented \textit{post hoc} for litigation---and must not rely on overbroad generalizations about the different talents, capacities, or preferences of men and women (\href{https://supreme.justia.com/cases/federal/us/518/515/}{id.}). Where a State unconstitutionally excludes women from a unique public educational opportunity, a remedial separate program that is unequal in both tangible resources and intangible prestige does not cure the violation (\href{https://supreme.justia.com/cases/federal/us/518/515/}{id. at 547--554}).

\textbf{Holding vs. Dicta:} The \textit{holding} is that Virginia failed to supply an exceedingly persuasive justification for categorically excluding qualified women from VMI, that the Fourth Circuit erred by applying deferential review and a ``substantive comparability'' standard to VWIL, and that VWIL failed to provide substantially equal educational opportunities under \textit{Sweatt v. Painter}, 339 U.S. 629 (1950) (\href{https://supreme.justia.com/cases/federal/us/518/515/}{518 U.S. at 534, 553--556}). By contrast, the Court's observations in Footnote 6 that strict scrutiny has ``thus far'' been reserved for race and national origin, its reassurance in Footnote 7 that it addressed ``specifically and only'' a ``unique'' educational opportunity rather than evenhanded public support for single-sex diversity generally, and its recognition that ``inherent differences'' between men and women may justify compensatory sex classifications that do not perpetuate women's inferiority are influential \textit{dicta} that framed subsequent debates without resolving them (\href{https://supreme.justia.com/cases/federal/us/518/515/}{id. at 532 n.6, 533--534 \& n.7}).

\section{Analysis of Opinions}

\subsection{Majority Opinion (Justice Ginsburg)}
Justice Ginsburg situated heightened scrutiny of sex classifications in the Nation's ``long and unfortunate history of sex discrimination,'' tracing the shift from deferential rational-basis cases such as \textit{Goesaert v. Cleary}, 335 U.S. 464 (1948), to the modern post-\textit{Reed v. Reed}, 404 U.S. 71 (1971), doctrine forbidding government from denying women ``full citizenship stature'' (\href{https://supreme.justia.com/cases/federal/us/518/515/}{518 U.S. at 531--532}). Without equating sex classifications to race for all purposes, she distilled the standard as ``skeptical scrutiny'' demanding an ``exceedingly persuasive justification'' that is genuine rather than \textit{post hoc} and free from overbroad gender stereotypes (\href{https://supreme.justia.com/cases/federal/us/518/515/}{id. at 531--534}). While acknowledging enduring physical differences between the sexes, she emphasized that sex classifications may be used to compensate women for economic disabilities or promote equal opportunity, but never ``to create or perpetuate the legal, social, and economic inferiority of women'' (\href{https://supreme.justia.com/cases/federal/us/518/515/}{id. at 533--534}).

Applying this framework to liability, Justice Ginsburg rejected both of Virginia's proffered justifications. First, the historical record refuted Virginia's claim that VMI was established or maintained to foster educational diversity: VMI was founded in 1839 when higher education was widely viewed as harmful to women, Virginia's other public single-sex schools had all become coeducational by the mid-1970s, the Commonwealth's 1990 higher-education commission explicitly called for nondiscrimination without regard to sex, and a plan affording a unique benefit solely to Virginia's sons while making no provision for its daughters did not reflect genuine diversity (\href{https://supreme.justia.com/cases/federal/us/518/515/}{id. at 535--540}). Second, Virginia's claim that admitting women would destroy VMI's adversative method rested on ``self-fulfilling prophec[ies]'' and average-tendency generalizations (\href{https://supreme.justia.com/cases/federal/us/518/515/}{id. at 540--545}). Because it was undisputed that some women wanted to attend VMI and could meet all of its physical and mental demands, and because federal service academies had successfully integrated women with manageable privacy and physical-training adjustments, excluding all women disregarded individual merit (\href{https://supreme.justia.com/cases/federal/us/518/515/}{id. at 540--546, 550--551 n.19}).

Turning to the remedy, Justice Ginsburg invoked \textit{Milliken v. Bradley}, 433 U.S. 267 (1977), and \textit{Sweatt v. Painter}, 339 U.S. 629 (1950), holding that a constitutional remedy must place victims in the position they would have occupied absent discrimination (\href{https://supreme.justia.com/cases/federal/us/518/515/}{518 U.S. at 547, 553--554}). VWIL failed completely: it denied women VMI's rigorous military training on the basis of generalizations about ``the way women are,'' and it was a ``pale shadow'' of VMI in curriculum, faculty credentials, funding, athletic facilities, prestige, and alumni influence (\href{https://supreme.justia.com/cases/federal/us/518/515/}{id. at 548--554}). Finally, the Fourth Circuit plainly erred by replacing heightened scrutiny with deferential review and its own ``substantive comparability'' test (\href{https://supreme.justia.com/cases/federal/us/518/515/}{id. at 554--556}).

\subsection{Concurrence in the Judgment (Chief Justice Rehnquist)}
Chief Justice Rehnquist agreed with both holdings---that VMI's male-only policy violated Equal Protection and that VWIL did not cure the violation---but objected to three aspects of the majority's analysis (\href{https://supreme.justia.com/cases/federal/us/518/515/}{518 U.S. at 558}). First, he criticized the Court's repeated elevation of ``exceedingly persuasive justification'' from a descriptive observation into the governing test, warning that it introduced unnecessary doctrinal uncertainty into the established \textit{Craig v. Boren}, 429 U.S. 190 (1976), intermediate-scrutiny standard (\href{https://supreme.justia.com/cases/federal/us/518/515/}{id. at 558--559}). Second, he refused to draw adverse inferences from Virginia's pre-1982 history, arguing that Virginia was not placed on notice of constitutional infirmity until \textit{Mississippi University for Women v. Hogan}, 458 U.S. 718 (1982); nevertheless, Virginia's failure after 1982 to provide a comparable public single-sex option for women rendered its single-sex diversity claim fatally one-sided, and the record lacked evidence that the adversative method itself was pedagogically superior (\href{https://supreme.justia.com/cases/federal/us/518/515/}{id. at 560--564}). Third, he defined the violation more narrowly---not as excluding women from VMI, but as maintaining an all-male public school without an equal-caliber all-female public school---meaning Virginia could have remedied the violation without coeducating VMI had it created a truly comparable, self-standing, adequately funded women's institution rather than the ``distinctly inferior'' VWIL program (\href{https://supreme.justia.com/cases/federal/us/518/515/}{id. at 565--566}).

\subsection{Dissenting Opinion (Justice Scalia)}
Justice Scalia mounted a sweeping originalist, institutional, and doctrinal critique. First, invoking tradition-based constitutional interpretation, he argued that when a practice not expressly barred by constitutional text enjoys an open, widespread, and unbroken tradition---like all-male military colleges and single-sex education---courts have no warrant to abolish it rather than leaving change to democratic processes (\href{https://www.law.cornell.edu/supct/html/94-1941.ZD.html}{518 U.S. at 568--570}). Second, he accused the majority of executing a ``de facto abandonment'' of intermediate scrutiny and covertly adopting strict scrutiny by requiring a perfect fit (invalidating VMI simply because ``some women'' could meet its standards) while playing ``Supreme Court peek-a-boo'' in Footnote 6 about whether strict scrutiny might formally apply (\href{https://www.law.cornell.edu/supct/html/94-1941.ZD.html}{id. at 571--574}). If the standard of review were reconsidered, he argued, rational-basis review would be more faithful to \textit{United States v. Carolene Products Co.}, 304 U.S. 144, 152--153 n.4 (1938), because women constitute a majority of the electorate rather than a ``discrete and insular minority'' (\href{https://www.law.cornell.edu/supct/html/94-1941.ZD.html}{id. at 574--576}). Third, applying traditional intermediate scrutiny, Justice Scalia contended that the majority improperly disregarded concurrent factual findings of both lower courts that single-sex education yields substantial pedagogical benefits, that Virginia's policy genuinely sought public-private diversity (including four state-assisted private women's colleges), and that admitting women would fundamentally alter VMI's adversative barracks culture (\href{https://www.law.cornell.edu/supct/html/94-1941.ZD.html}{id. at 576--592}). Finally, he warned that the Court's rationale rendered public single-sex education ``functionally dead'' and threatened public financial support and tax-exempt status for private single-sex colleges (\href{https://www.law.cornell.edu/supct/html/94-1941.ZD.html}{id. at 595--600}).

\section{Examples: Future Applications}
\begin{enumerate}
\item \textbf{Same-Side (Unconstitutional): State Maritime Academy Excluding Women.} A state-funded maritime academy restricts admission to men, arguing that sea-service leadership training relies on harsh shipboard discipline that most women would not prefer and that women may attend coeducational state engineering schools instead. Under \textit{Virginia}, this categorical exclusion fails heightened scrutiny because the State relies on overbroad generalizations about average preferences to bar qualified individual women from a unique, prestigious state-funded credential.
\item \textbf{Same-Side (Unconstitutional): Separate, Underfunded Public STEM Academy for Girls.} After a court invalidates an elite all-boys public STEM magnet high school, the school district keeps the boys' school intact and opens an all-girls ``STEM Leadership Wing'' inside a general high school with fewer AP science courses, no robotics lab, and an emphasis on ``collaborative communication.'' Under \textit{Virginia}'s remedial holding (and \textit{Sweatt}), this ``pale shadow'' remedy fails to place excluded girls in the position they would have occupied absent discrimination.
\item \textbf{Opposite-Side (Constitutional): Remedial Single-Sex Scholarship for Underrepresented Women in Engineering.} A state university establishes a targeted fellowship program for women entering mechanical engineering, backed by contemporaneous data showing women comprise only 14\% of graduates and face documented historical barriers in the field, while leaving all general university scholarships open to men. Under \textit{Virginia}, sex classifications used to compensate women for documented economic or educational disabilities and to ``advance full development of the talent and capacities of our Nation's people''---rather than to perpetuate inferiority---can satisfy intermediate scrutiny (\href{https://supreme.justia.com/cases/federal/us/518/515/}{518 U.S. at 533--534}).
\item \textbf{Opposite-Side (Constitutional): Sex-Separated Barracks Showers and Physiological Physical-Fitness Benchmarks at a Coeducational Police Academy.} A state police academy admits both men and women to identical tactical and academic training, but provides sex-separated sleeping and shower facilities for privacy and adjusts upper-body strength test benchmarks to reflect average physiological differences while maintaining an equally rigorous standard of job-related fitness. Justice Ginsburg explicitly recognized in Footnote 19 that privacy accommodations and minimum essential physical-training adjustments based on enduring physical differences are compatible with equal protection (\href{https://supreme.justia.com/cases/federal/us/518/515/}{id. at 533, 550--551 n.19}).
\item \textbf{Fence-Sitter (Close Case): Simultaneous Creation of Twin Public Single-Sex Leadership High Schools.} A large urban school district simultaneously opens two brand-new, voluntary public high schools---one for boys and one for girls---with identical per-pupil funding, identical STEM and humanities curricula, comparable facilities, equally credentialed faculties, and a documented pedagogical goal of improving graduation rates, while maintaining high-performing coeducational options. Because \textit{Virginia} expressly reserved the question of evenhanded, truly equal single-sex offerings in Footnote 7 (and cited Judge Phillips's ``simultaneous creation'' paradigm), this program sits directly on the doctrinal fault line: defenders would argue it avoids the stigma and disparity of VMI/VWIL, whereas challengers would argue that sex segregation itself rests on impermissible generalizations about how boys and girls learn (\href{https://supreme.justia.com/cases/federal/us/518/515/}{id. at 534 n.7, 547 n.17}).
\end{enumerate}

\section{Critique}
\textbf{Originalist and Textualist Critiques:} Conservative and originalist scholars echo Justice Scalia's objection that the Fourteenth Amendment's original public meaning in 1868 did not forbid sex-segregated public colleges, and that \textit{Virginia} destabilized equal-protection tiers by turning intermediate scrutiny into ``strict scrutiny in disguise'' without candidly adopting strict scrutiny (\href{https://www.law.cornell.edu/supct/html/94-1941.ZD.html}{518 U.S. at 571--574}). Recent scholarship examining \href{https://www.supremecourt.gov/opinions/24pdf/23-477_2cp3.pdf}{\textit{United States v. Skrmetti}} (2025) notes that Scalia's \textit{VMI} dissent remains the ``locus classicus'' of the conservative constitutional critique of sex-equality heightened scrutiny, though Justice Amy Coney Barrett's \textit{Skrmetti} concurrence explicitly reaffirmed intermediate scrutiny for sex classifications and rejected Scalia's \textit{Carolene Products} political-power objection (\href{https://direct.mit.edu/ajle/article/doi/10.1162/AJLE.a.107/139252/AMY-CONEY-BARRETT-S-SEX-EQUALITY-CONSTITUTION-The}{AJLE 2026}).

\textbf{Progressive, Feminist, and Doctrinal Perspectives:} Many progressive and doctrinal scholars praised \textit{Virginia} as a landmark articulation of anti-subordination principles. Cass Sunstein described Justice Ginsburg's opinion as ``ambitious'' and ``nonminimalist,'' articulating a substantive vision of sex equality that acknowledges biological differences while prohibiting the state from converting difference into social or economic hierarchy (\href{https://ir.law.fsu.edu/cgi/viewcontent.cgi?article=1533&context=articles}{FSU Law Review, citing Sunstein, 110 Harv. L. Rev. 4, 74 (1996)}). Mary Anne Case argued that Justice Scalia accurately described \textit{VMI}'s demand for near-perfect proxies when generalizations exclude capable individuals, but contended that this demand was entirely consistent with the Court's post-\textit{Reed} sex-discrimination precedent rather than a departure (\href{https://chicagounbound.uchicago.edu/cgi/viewcontent.cgi?article=2109&context=journal_articles}{Case, Chicago Unbound}). Conversely, some feminist critics, including Katharine Bartlett, faulted the litigation and opinion for accepting VMI's hyper-masculine ``adversative method'' as a valuable baseline rather than interrogating how the institution's pedagogy itself demeaned and subordinated women (\href{https://scholarship.law.duke.edu/cgi/viewcontent.cgi?article=2936&context=faculty_scholarship}{Bartlett, Duke Law}).

\textbf{Underlying Competing Values:} At bottom, \textit{Virginia} exposes the tension between \textit{individualized equal opportunity} (anti-stereotyping and anti-subordination) and \textit{institutional pluralism and democratic federalism} (allowing states to preserve longstanding traditions and experiment with diverse educational models tailored to statistical group tendencies).

\section{Key Quotations}
\begin{quote}
``Parties who seek to defend gender-based government action must demonstrate an `exceedingly persuasive justification' for that action.''
\end{quote}
--- Justice Ruth Bader Ginsburg, Opinion of the Court, 518 U.S. at 531 (\href{https://supreme.justia.com/cases/federal/us/518/515/}{Justia}; \href{https://www.courtlistener.com/opinion/118056/united-states-v-virginia/}{CourtListener}).

\begin{quote}
``The justification must be genuine, not hypothesized or invented post hoc in response to litigation. And it must not rely on overbroad generalizations about the different talents, capacities, or preferences of males and females.''
\end{quote}
--- Justice Ruth Bader Ginsburg, Opinion of the Court, 518 U.S. at 533 (\href{https://supreme.justia.com/cases/federal/us/518/515/}{Justia}).

\begin{quote}
````Inherent differences'' between men and women, we have come to appreciate, remain cause for celebration, but not for denigration of the members of either sex or for artificial constraints on an individual's opportunity.''
\end{quote}
--- Justice Ruth Bader Ginsburg, Opinion of the Court, 518 U.S. at 533 (\href{https://supreme.justia.com/cases/federal/us/518/515/}{Justia}).

\begin{quote}
``While terms like `important governmental objective' and `substantially related' are hardly models of precision, they have more content and specificity than does the phrase `exceedingly persuasive justification.' That phrase is best confined, as it was first used, as an observation on the difficulty of meeting the applicable test, not as a formulation of the test itself.''
\end{quote}
--- Chief Justice William H. Rehnquist, Concurring in the Judgment, 518 U.S. at 559 (\href{https://supreme.justia.com/cases/federal/us/518/515/}{Justia}).

\begin{quote}
``The virtue of a democratic system with a First Amendment is that it readily enables the people, over time, to be persuaded that what they took for granted is not so, and to change their laws accordingly. That system is destroyed if the smug assurances of each age are removed from the democratic process and written into the Constitution.''
\end{quote}
--- Justice Antonin Scalia, Dissenting, 518 U.S. at 567 (\href{https://www.law.cornell.edu/supct/html/94-1941.ZD.html}{Cornell LII}).

\end{document}
